\chapter{Conformal Killing Vectors and the Conformal Group}
\label{ch:volume-iii-conformal-killing-vectors}

\section{The Conformal Killing Equation}

Work first in flat Euclidean space \((\R^D,\delta_{\mu\nu})\) with
\(D\ge3\).  A smooth vector field \(\epsilon^\mu(x)\) generates an
infinitesimal coordinate transformation
\[
  x^\mu\longmapsto x^\mu+\epsilon^\mu(x).
\]
It is conformal when the change of the metric is proportional to the metric:
\[
  \mathcal L_\epsilon\delta_{\mu\nu}
  =
  \partial_\mu\epsilon_\nu+\partial_\nu\epsilon_\mu
  =
  2\sigma_\epsilon(x)\delta_{\mu\nu}
\]
for some scalar function \(\sigma_\epsilon\).  Taking the trace gives
\[
  \sigma_\epsilon(x)
  =
  \frac1D\partial_\rho\epsilon^\rho(x).
\]
Thus the conformal Killing equation is
\begin{equation}
  \partial_\mu\epsilon_\nu+\partial_\nu\epsilon_\mu
  -
  \frac2D\delta_{\mu\nu}\partial_\rho\epsilon^\rho
  =
  0.
  \label{eq:conformal-killing-equation}
\end{equation}

If \(T_{\mu\nu}\) is symmetric, conserved, and traceless, then
\[
  J^\mu_\epsilon(x)=T^{\mu}{}_{\nu}(x)\epsilon^\nu(x)
\]
is conserved:
\[
  \partial_\mu J^\mu_\epsilon
  =
  (\partial_\mu T^{\mu}{}_{\nu})\epsilon^\nu
  +
  T^{\mu\nu}\partial_\mu\epsilon_\nu
  =
  \frac12T^{\mu\nu}
  (\partial_\mu\epsilon_\nu+\partial_\nu\epsilon_\mu)
  =
  \sigma_\epsilon T^\mu{}_\mu
  =
  0.
\]
The corresponding charge on a closed codimension-one surface \(\Sigma\) is
\[
  Q_\epsilon(\Sigma)
  =
  \int_\Sigma \dd\sigma\,n_\mu J^\mu_\epsilon,
\]
where \(n_\mu\) is the unit normal.  Conservation makes the charge invariant
under deformations of \(\Sigma\) that do not cross operator insertions.

\section{Solutions in Dimension at Least Three}

For \(D\ge3\), differentiating
\eqref{eq:conformal-killing-equation} and combining permutations of the
indices gives
\[
  \partial_\mu\partial_\nu\epsilon_\rho
  =
  \delta_{\rho\nu}\partial_\mu\sigma_\epsilon
  +
  \delta_{\rho\mu}\partial_\nu\sigma_\epsilon
  -
  \delta_{\mu\nu}\partial_\rho\sigma_\epsilon .
\]
Taking another derivative and antisymmetrizing shows
\[
  \partial_\mu\partial_\nu\partial_\rho\sigma_\epsilon=0.
\]
Hence \(\sigma_\epsilon\) is at most linear in \(x\), and
\(\epsilon^\mu\) is at most quadratic.  The general solution is
\begin{equation}
  \epsilon^\mu(x)
  =
  a^\mu
  +
  \omega^\mu{}_\nu x^\nu
  +
  \lambda x^\mu
  +
  2(b\cdot x)x^\mu-b^\mu x^2,
  \label{eq:ckv-general-solution}
\end{equation}
where \(a^\mu\), \(b^\mu\), and \(\lambda\) are constants and
\(\omega_{\mu\nu}=-\omega_{\nu\mu}\).  These terms generate translations,
rotations, dilatations, and special conformal transformations, respectively.

The associated differential operators are
\begin{equation}
  P_\mu=\partial_\mu,\qquad
  M_{\mu\nu}=x_\mu\partial_\nu-x_\nu\partial_\mu,
  \qquad
  D=x^\mu\partial_\mu,
  \label{eq:conformal-differential-generators-1}
\end{equation}
and
\begin{equation}
  K_\mu=x^2\partial_\mu-2x_\mu x^\rho\partial_\rho .
  \label{eq:conformal-differential-generators-2}
\end{equation}
The signs in \eqref{eq:conformal-differential-generators-2} are chosen so
that the finite special conformal transformation below has parameter \(b^\mu\)
with denominator \(1-2b\cdot x+b^2x^2\).

\begin{figure}[H]
  \centering
  \resizebox{0.96\textwidth}{!}{%
  \begin{tikzpicture}[x=1cm,y=1cm,every node/.style={font=\scriptsize}]
    \node[font=\small] at (0,2.25) {conformal Killing vector};
    \draw[rounded corners=4pt,thick] (-2.0,0.45) rectangle (2.0,1.55);
    \node at (0,1.00) {\(\epsilon^\mu(x)\)};

    \draw[-{Latex[length=2mm]},thick] (2.25,1.0)--(3.45,1.0);

    \begin{scope}[shift={(5.35,0)}]
      \node[font=\small] at (0,2.25) {finite-dimensional data};
      \draw[rounded corners=4pt,thick] (-2.55,0.35) rectangle (2.55,1.65);
      \node at (0,1.27) {\(a^\mu,\ \omega_{\mu\nu},\ \lambda,\ b^\mu\)};
      \node[align=center] at (0,0.72)
        {translation, rotation,\\dilatation, special conformal};
    \end{scope}

    \draw[-{Latex[length=2mm]},thick] (7.95,1.0)--(9.15,1.0);

    \begin{scope}[shift={(11.0,0)}]
      \node[font=\small] at (0,2.25) {charges};
      \draw[rounded corners=4pt,thick] (-1.75,0.45) rectangle (1.75,1.55);
      \node at (0,1.00) {\(Q_\epsilon=\int_\Sigma n_\mu T^\mu{}_\nu\epsilon^\nu\)};
    \end{scope}
  \end{tikzpicture}
  }
  \caption{In \(D\ge3\), a conformal Killing vector is determined by finitely
  many constants.  The stress tensor converts such a vector field into a
  conserved charge.}
  \label{fig:ckv-data-to-charge}
\end{figure}

\section{Finite Transformations}

A finite conformal transformation is a diffeomorphism
\[
  f:\R^D\supset U\longrightarrow f(U)\subset\R^D
\]
whose Jacobian satisfies
\begin{equation}
  \frac{\partial f^\rho}{\partial x^\mu}
  \frac{\partial f^\sigma}{\partial x^\nu}
  \delta_{\rho\sigma}
  =
  \Omega_f(x)^2\delta_{\mu\nu},
  \qquad
  \Omega_f(x)>0.
  \label{eq:finite-conformal-map}
\end{equation}
Translations, rotations, and dilatations are immediate examples.  A further
basic example is inversion,
\begin{equation}
  I(x)^\mu=\frac{x^\mu}{x^2},
  \qquad
  x^2=\delta_{\mu\nu}x^\mu x^\nu .
  \label{eq:inversion-map}
\end{equation}
Its Jacobian is
\[
  \frac{\partial I^\mu}{\partial x^\nu}
  =
  \frac1{x^2}
  \left(
  \delta^\mu{}_\nu-2\frac{x^\mu x_\nu}{x^2}
  \right),
\]
and the matrix in parentheses is orthogonal.  Therefore
\[
  \Omega_I(x)=\frac1{x^2}.
\]

The composition
\[
  I\circ T_{-b}\circ I
\]
with translation \(T_{-b}(x)=x-b\) gives a special conformal
transformation
\begin{equation}
  x^\mu
  \longmapsto
  x'^\mu
  =
  \frac{x^\mu-b^\mu x^2}{1-2b\cdot x+b^2x^2}.
  \label{eq:finite-special-conformal}
\end{equation}
Expanding to first order in \(b\) gives
\[
  \delta x^\mu
  =
  2(b\cdot x)x^\mu-b^\mu x^2,
\]
which is the last term in \eqref{eq:ckv-general-solution}.

\section{The Conformal Algebra}

Let the quantum charges associated with the vector fields
\eqref{eq:conformal-differential-generators-1} and
\eqref{eq:conformal-differential-generators-2} be denoted by
\[
  \widehat P_\mu,\qquad
  \widehat J_{\mu\nu},\qquad
  \widehat D,\qquad
  \widehat K_\mu .
\]
Their commutators are the charge realization of the Lie bracket of conformal
Killing vector fields.  In Lorentzian signature, with metric
\(\eta_{\mu\nu}\), the algebra is
\begin{align}
  [\widehat D,\widehat P_\mu]&=\ii\widehat P_\mu,
  &
  [\widehat D,\widehat K_\mu]&=-\ii\widehat K_\mu,
  &
  [\widehat D,\widehat J_{\mu\nu}]&=0,
  \label{eq:conf-alg-dilatation}\\
  [\widehat P_\mu,\widehat K_\nu]
  &=
  2\ii(\eta_{\mu\nu}\widehat D-\widehat J_{\mu\nu}),
  \label{eq:conf-alg-pk}
\end{align}
together with the Poincare commutators
\begin{align}
  [\widehat J_{\mu\nu},\widehat P_\rho]
  &=
  \ii(\eta_{\rho\mu}\widehat P_\nu-\eta_{\rho\nu}\widehat P_\mu),
  \\
  [\widehat J_{\mu\nu},\widehat J_{\rho\sigma}]
  &=
  \ii(
  \eta_{\mu\rho}\widehat J_{\nu\sigma}
  -\eta_{\nu\rho}\widehat J_{\mu\sigma}
  -\eta_{\mu\sigma}\widehat J_{\nu\rho}
  +\eta_{\nu\sigma}\widehat J_{\mu\rho}).
\end{align}
All unlisted commutators among the \(\widehat P_\mu\)'s and among the
\(\widehat K_\mu\)'s vanish.

The algebra is isomorphic to \(\mathfrak{so}(D,2)\).  Introduce indices
\[
  M,N=-1,0,1,\ldots,D
\]
and a metric
\[
  \eta_{MN}^{(D,2)}=\operatorname{diag}(-1,-1,+1,\ldots,+1).
\]
The generators \(\widehat{\mathcal J}_{MN}=-\widehat{\mathcal J}_{NM}\) are
related to the conformal generators by
\[
  \widehat{\mathcal J}_{\mu\nu}=\widehat J_{\mu\nu},\qquad
  \widehat{\mathcal J}_{D,-1}=\widehat D,\qquad
  \widehat{\mathcal J}_{D,\mu}=\frac12(\widehat P_\mu-\widehat K_\mu),
  \qquad
  \widehat{\mathcal J}_{-1,\mu}=\frac12(\widehat P_\mu+\widehat K_\mu).
\]
Then
\[
  [\widehat{\mathcal J}_{MN},\widehat{\mathcal J}_{PQ}]
  =
  \ii(
  \eta_{MP}^{(D,2)}\widehat{\mathcal J}_{NQ}
  -\eta_{NP}^{(D,2)}\widehat{\mathcal J}_{MQ}
  -\eta_{MQ}^{(D,2)}\widehat{\mathcal J}_{NP}
  +\eta_{NQ}^{(D,2)}\widehat{\mathcal J}_{MP}).
\]

\section{Two Euclidean Dimensions}

For \(D=2\), the local conformal Killing equation has infinitely many
solutions.  In complex coordinates
\[
  z=x^1+\ii x^2,\qquad
  \bar z=x^1-\ii x^2,
\]
the local conformal transformations are
\[
  z\mapsto f(z),\qquad
  \bar z\mapsto \bar f(\bar z)
\]
in Euclidean signature, with \(f\) holomorphic in its domain.  The global
conformal transformations of the Riemann sphere are the Möbius maps
\[
  z\mapsto \frac{az+b}{cz+d},
  \qquad
  ad-bc\neq0.
\]
The infinite-dimensional enhancement is the entry point to the Virasoro
algebra, treated later after the higher-dimensional conformal representation
theory has been established.
